\section{Mata Kuliah dalam Kurikulum dan RPS Berbasis OBE}
\label{sec:ch01-kurikulum}

Mata kuliah Teknik Kompilasi (MWCA017) diselenggarakan pada Semester
Ganjil dengan bobot 3 SKS. Mata kuliah ini mendukung dua Capaian
Pembelajaran Lulusan (CPL), yaitu CPL 02 tentang cara kerja sistem
komputer dan CPL 03 tentang perancangan serta analisis algoritma.
Turunan kedua CPL tersebut dinyatakan dalam empat CPMK.

\begin{table}[H]
\centering
\caption{Keterkaitan CPL, CPMK, dan fokus buku}
\label{tab:ch01-cpl-cpmk}
\begin{tabularx}{\textwidth}{@{}lX@{}}
\toprule
\textbf{Komponen} & \textbf{Rumusan dan fokus} \\
\midrule
CPL 02 & Memahami cara kerja sistem komputer serta mampu merancang
dan mengembangkan algoritma untuk memecahkan masalah. \\
CPL 03 & Mampu merancang dan menganalisis algoritma secara optimal
serta menerapkannya pada bahasa pemrograman tertentu. \\
\midrule
CPMK 01 & Menjelaskan tahapan kompilasi: leksikal, sintaksis, dan
semantik. \\
CPMK 02 & Memahami konsep otomata hingga. \\
CPMK 03 & Merancang tata bahasa bebas konteks untuk bahasa sederhana. \\
CPMK 04 & Mengimplementasikan parser sederhana untuk validasi kode. \\
\bottomrule
\end{tabularx}
\end{table}

RPS membagi 16 pertemuan menjadi 14 pertemuan materi, satu Ujian Tengah
Semester (UTS), dan satu Ujian Akhir Semester (UAS). Buku ini memetakan
keempat belas topik tersebut ke empat belas bab inti, sehingga setiap
pertemuan memiliki rujukan bab yang jelas.

\begin{warningbox}
Buku ini menggunakan data identitas dari ketentuan program studi:
Universitas Bale Bandung, Fakultas Teknologi Informasi, kode MWCA017,
dan Semester Ganjil. Data institusi pada berkas RPS asli merupakan
placeholder dan tidak dipakai.
\end{warningbox}
